% GENERATED FILE. Edit canonical lessons, then run npm run generate:books.
% canonical-source-hash: fnv1a64:d1b58d77c2fbeac4
% canonical-lessons: TE-C526-niyamincu, TE-C526-nirvahincu, TE-C526-paryaveksincu, TE-C526-sahakarincu, TE-C526-siksana

\chapter{Running a Team}
\label{ch:te-running-a-team}
\hlchaptermodality{526}

\begin{chapteropening}
\textbf{By the end of this chapter:} \emph{I can say that someone appointed, managed, supervised or trained people.}

Complete the last lesson of chapter 526: I can say that someone appointed, managed, supervised or trained people.
\end{chapteropening}

\section[\texorpdfstring{\te{నియమించు} (niyamiñcu)}{నియమించు (niyamiñcu)}]{\te{నియమించు} (niyamiñcu) — to appoint; to assign}

\label{lesson:TE-C526-niyamincu}



\begin{quote}
Before the new one: say the Telugu for a wage, then the Telugu for to toil.
\end{quote}

\subsection*{You'll want to know: \te{నియమించు}}
\textbf{\te{నియమించు}} — \emph{niyamiñcu} — \textquotedblleft{}to appoint; to assign\textquotedblright{}.

\begin{quote}
\te{కంపెనీ} \te{కొత్త} \te{మేనేజరును} \te{నియమించింది}.
\end{quote}

\begin{quote}
\emph{kampenī kotta mēnējarunu niyamiñcindi.}
\end{quote}

\begin{quote}
The company appointed a new manager.
\end{quote}

Past (I …): \te{నియమించాను}, \emph{niyamiñcānu}.

-\te{గా} \te{నియమించు}, to appoint as: \te{ఆమెను} \te{మేనేజరుగా} \te{నియమించారు}.

\begin{cousinweb}
\textbf{\te{నియమించు}} (\emph{niyamiñcu}) is built on Sanskrit \emph{niyama}, \textquotedblleft{}restraint; a fixed rule\textquotedblright{} (CDIAL 7250), with the Telugu verb ending -\te{ఇంచు}: to fix someone in a place.
\end{cousinweb}

\begin{grammarlens}[title={What you've built}]
One more word for describing what you lived and felt.
\end{grammarlens}

\subsection*{Your turn}
\begin{itemize}
\raggedright
  \item \emph{Say it:} \emph{niyamiñcu}
  \item \emph{Say it:} \emph{niyamiñcu}, once more
  \item \emph{Say it:} \emph{kampenī kotta mēnējarunu niyamiñcindi.} — The company appointed a new manager.
  \item \emph{Recall:} say the Telugu for a wage, then the Telugu for to toil, then say \emph{niyamiñcu} again
\end{itemize}

\begin{tcolorbox}[breakable,colback=teal!4,colframe=teal!35!black,title={Before you move on}]
What does \te{నియమించు} mean? (\textquotedblleft{}To appoint; to assign\textquotedblright{}.) Say it once more.
\end{tcolorbox}
\section[\texorpdfstring{\te{నిర్వహించు} (nirvahiñcu)}{నిర్వహించు (nirvahiñcu)}]{\te{నిర్వహించు} (nirvahiñcu) — to manage; to run}

\label{lesson:TE-C526-nirvahincu}



\begin{quote}
Before the new one: say the Telugu for to toil, then the Telugu for to appoint.
\end{quote}

\subsection*{You'll want to know: \te{నిర్వహించు}}
\textbf{\te{నిర్వహించు}} — \emph{nirvahiñcu} — \textquotedblleft{}to manage; to run\textquotedblright{}.

\begin{quote}
\te{మా} \te{అక్క} \te{ఒక} \te{చిన్న} \te{దుకాణం} \te{నిర్వహిస్తుంది}.
\end{quote}

\begin{quote}
\emph{mā akka oka cinna dukāṇaṁ nirvahistundi.}
\end{quote}

\begin{quote}
My elder sister runs a small shop.
\end{quote}

Past (I …): \te{నిర్వహించాను}, \emph{nirvahiñcānu}.

Also to hold an event: \te{సమావేశం} \te{నిర్వహించు}, to hold a meeting.

\begin{cousinweb}
\textbf{\te{నిర్వహించు}} (\emph{nirvahiñcu}) is built on Sanskrit \emph{nirvahaṇa}, \textquotedblleft{}carrying through, completion\textquotedblright{} (Monier-Williams), with the Telugu verb ending -\te{ఇంచు}.
\end{cousinweb}

\begin{grammarlens}[title={What you've built}]
One more word for describing what you lived and felt.
\end{grammarlens}

\subsection*{Your turn}
\begin{itemize}
\raggedright
  \item \emph{Say it:} \emph{nirvahiñcu}
  \item \emph{Say it:} \emph{nirvahiñcu}, once more
  \item \emph{Say it:} \emph{mā akka oka cinna dukāṇaṁ nirvahistundi.} — My elder sister runs a small shop.
  \item \emph{Recall:} say the Telugu for to toil, then the Telugu for to appoint, then say \emph{nirvahiñcu} again
\end{itemize}

\begin{tcolorbox}[breakable,colback=teal!4,colframe=teal!35!black,title={Before you move on}]
What does \te{నిర్వహించు} mean? (\textquotedblleft{}To manage; to run\textquotedblright{}.) Say it once more.
\end{tcolorbox}
\section[\texorpdfstring{\te{పర్యవేక్షించు} (paryavēkṣiñcu)}{పర్యవేక్షించు (paryavēkṣiñcu)}]{\te{పర్యవేక్షించు} (paryavēkṣiñcu) — to supervise; to oversee}

\label{lesson:TE-C526-paryaveksincu}



\begin{quote}
Before the new one: say the Telugu for to appoint, then the Telugu for to manage.
\end{quote}

\subsection*{You'll want to know: \te{పర్యవేక్షించు}}
\textbf{\te{పర్యవేక్షించు}} — \emph{paryavēkṣiñcu} — \textquotedblleft{}to supervise; to oversee\textquotedblright{}.

\begin{quote}
\te{ఇంజనీరు} \te{పనిని} \te{దగ్గరుండి} \te{పర్యవేక్షించారు}.
\end{quote}

\begin{quote}
\emph{iñjanīru panini daggaruṇḍi paryavēkṣiñcāru.}
\end{quote}

\begin{quote}
The engineer oversaw the work in person.
\end{quote}

Past (I …): \te{పర్యవేక్షించాను}, \emph{paryavēkṣiñcānu}.

The noun is \te{పర్యవేక్షణ}, supervision.

\begin{cousinweb}
\textbf{\te{పర్యవేక్షించు}} (\emph{paryavēkṣiñcu}) is a learned formation built from Sanskrit \emph{pari-}, \textquotedblleft{}around\textquotedblright{}, and \emph{avekṣaṇa}, \textquotedblleft{}looking at\textquotedblright{} (Monier-Williams), with the Telugu verb ending -\te{ఇంచు}: to look all round a task.
\end{cousinweb}

\begin{grammarlens}[title={What you've built}]
One more word for describing what you lived and felt.
\end{grammarlens}

\subsection*{Your turn}
\begin{itemize}
\raggedright
  \item \emph{Say it:} \emph{paryavēkṣiñcu}
  \item \emph{Say it:} \emph{paryavēkṣiñcu}, once more
  \item \emph{Say it:} \emph{iñjanīru panini daggaruṇḍi paryavēkṣiñcāru.} — The engineer oversaw the work in person.
  \item \emph{Recall:} say the Telugu for to appoint, then the Telugu for to manage, then say \emph{paryavēkṣiñcu} again
\end{itemize}

\begin{tcolorbox}[breakable,colback=teal!4,colframe=teal!35!black,title={Before you move on}]
What does \te{పర్యవేక్షించు} mean? (\textquotedblleft{}To supervise; to oversee\textquotedblright{}.) Say it once more.
\end{tcolorbox}
\section[\texorpdfstring{\te{సహకరించు} (sahakariñcu)}{సహకరించు (sahakariñcu)}]{\te{సహకరించు} (sahakariñcu) — to cooperate; to help}

\label{lesson:TE-C526-sahakarincu}



\begin{quote}
Before the new one: say the Telugu for to manage, then the Telugu for to supervise.
\end{quote}

\subsection*{You'll want to know: \te{సహకరించు}}
\textbf{\te{సహకరించు}} — \emph{sahakariñcu} — \textquotedblleft{}to cooperate; to help\textquotedblright{}.

\begin{quote}
\te{అందరూ} \te{సహకరిస్తే} \te{పని} \te{త్వరగా} \te{అవుతుంది}.
\end{quote}

\begin{quote}
\emph{andarū sahakaristē pani tvaragā avutundi.}
\end{quote}

\begin{quote}
If everyone cooperates, the work gets done quickly.
\end{quote}

Past (I …): \te{సహకరించాను}, \emph{sahakariñcānu}.

-\te{కి} \te{సహకరించు}, to cooperate with someone.

\begin{cousinweb}
\textbf{\te{సహకరించు}} (\emph{sahakariñcu}) is built on Sanskrit \emph{sahakāra}, \textquotedblleft{}assistance\textquotedblright{} (CDIAL 13299), from \emph{saha}, \textquotedblleft{}with\textquotedblright{}, and \emph{kāra}, \textquotedblleft{}doing\textquotedblright{}.
\end{cousinweb}

\begin{grammarlens}[title={What you've built}]
One more word for describing what you lived and felt.
\end{grammarlens}

\subsection*{Your turn}
\begin{itemize}
\raggedright
  \item \emph{Say it:} \emph{sahakariñcu}
  \item \emph{Say it:} \emph{sahakariñcu}, once more
  \item \emph{Say it:} \emph{andarū sahakaristē pani tvaragā avutundi.} — If everyone cooperates, the work gets done quickly.
  \item \emph{Recall:} say the Telugu for to manage, then the Telugu for to supervise, then say \emph{sahakariñcu} again
\end{itemize}

\begin{tcolorbox}[breakable,colback=teal!4,colframe=teal!35!black,title={Before you move on}]
What does \te{సహకరించు} mean? (\textquotedblleft{}To cooperate; to help\textquotedblright{}.) Say it once more.
\end{tcolorbox}
\section[\texorpdfstring{\te{శిక్షణ} (śikṣaṇa)}{శిక్షణ (śikṣaṇa)}]{\te{శిక్షణ} (śikṣaṇa) — training}

\label{lesson:TE-C526-siksana}



\begin{quote}
Before the new one: say the Telugu for to supervise, then the Telugu for to cooperate.
\end{quote}

\subsection*{You'll want to know: \te{శిక్షణ}}
\textbf{\te{శిక్షణ}} — \emph{śikṣaṇa} — \textquotedblleft{}training\textquotedblright{}.

\begin{quote}
\te{కొత్త} \te{ఉద్యోగులకు} \te{ఒక} \te{నెల} \te{శిక్షణ} \te{ఇచ్చారు}.
\end{quote}

\begin{quote}
\emph{kotta udyōgulaku oka nela śikṣaṇa iccāru.}
\end{quote}

\begin{quote}
The new employees were given a month of training.
\end{quote}

\te{శిక్షణ} \te{ఇచ్చు}, to train someone; \te{శిక్షణ} \te{పొందు}, to be trained. Not \te{శిక్ష}, a punishment.

\begin{cousinweb}
\textbf{\te{శిక్షణ}} (\emph{śikṣaṇa}) is borrowed from Sanskrit \emph{śikṣaṇa}, \textquotedblleft{}the act of learning; teaching\textquotedblright{} (CDIAL 12429).
\end{cousinweb}

\begin{grammarlens}[title={What you've built}]
One more word for describing what you lived and felt.
\end{grammarlens}

\subsection*{Your turn}
\begin{itemize}
\raggedright
  \item \emph{Say it:} \emph{śikṣaṇa}
  \item \emph{Say it:} \emph{śikṣaṇa}, once more
  \item \emph{Say it:} \emph{kotta udyōgulaku oka nela śikṣaṇa iccāru.} — The new employees were given a month of training.
  \item \emph{Recall:} say the Telugu for to supervise, then the Telugu for to cooperate, then say \emph{śikṣaṇa} again
\end{itemize}

\begin{tcolorbox}[breakable,colback=teal!4,colframe=teal!35!black,title={Before you move on}]
What does \te{శిక్షణ} mean? (\textquotedblleft{}Training\textquotedblright{}.) Say it once more.
\end{tcolorbox}
